% Embeddable research section. Requires amsmath, amssymb, amsthm, booktabs,
% array, and hyperref. Bibliography entries are in NONINTEGRAL-bib.tex.
\section{Dimension characters, central restrictions, and a quaternionic model}
\label{sec:r11-nonintegral}

Throughout this section, $R_n$ abbreviates the supplied based ring
\texttt{R11\_1\_}$n$. The integer $n$ is its physical line number in the
rank-eleven multiplicity-one input, and basis indices run from $0$ to $10$,
with $0$ the unit. All statements concern the entire based multiplication
tensor. In particular, an obstruction to a complex categorification is
distinguished from an obstruction that initially assumes a spherical or
pseudo-unitary structure.

The results below exclude thirty-eight rings as Grothendieck rings of
complex fusion categories. They also construct a unitary categorification
of $R_{179}$, of dimension $72$. The categorical ingredients are established
results; the uniform applications, explicit witnesses, and finite input
identifications are recorded here with exact replay certificates. These
results do not assert a classification of associators on any realizable
based ring.

\subsection{Exact passage from arbitrary categorifications to spherical ones}
\label{subsec:r11-signed-bridge}

We first explain the scope of the computational reduction used below.
Sphericalize a hypothetical complex fusion category $\mathcal C$, and
choose a lift of each simple object, with the unit and duality normalized.
Write $\varepsilon$ for the central sign object. In multiplicity one, the
quotient of the sphericalization's Grothendieck ring by
$[\varepsilon]+1$ has structure constants
\[
 T_{ij}^{k}=\sigma_{ij}^{k}N_{ij}^{k},\qquad
 \sigma_{ij}^{k}\in\{1,-1\}
 \quad\text{when }N_{ij}^{k}=1.
\]
These constants satisfy associativity, the unit identities, cyclic trace
symmetry, and reversal with duals. The trace triple $(i,j,k^*)$ is identified
only under cyclic rotation and
$(a,b,c)\mapsto(c^*,b^*,a^*)$; arbitrary permutations, or commutativity of
the signed quotient, are not imposed. Changing lifts replaces
$T_{ij}^{k}$ by $s_i s_j s_k T_{ij}^{k}$, where
$s_0=1$ and $s_i=s_{i^*}\in\{1,-1\}$.

\paragraph{Completeness of the sign calculation.}
Put $\sigma=(-1)^x$ on each oriented orbit. Every associativity identity
becomes an exact equation
\[
 \sum_m c_m(-1)^{\langle m,x\rangle}=0.
\]
If one coefficient has absolute value equal to the sum of the absolute
values of the others, the positive solution and equality in the triangle
inequality force the participating sign monomials to agree. Each resulting
linear relation over $\mathbb F_2$ is replayed from its full associativity
identity. An independently checked gauge matrix proves that the chosen
normalization meets every gauge orbit. All remaining assignments are then
enumerated and checked.

For the fifty-nine inputs in the recorded nonintegral audit, at most one
binary variable remains after this reduction. The complete domain has
seventy-one assignments: fifty-nine positive original tables and twelve
alternative associative signed tables. Every alternative is excluded by
its complete scalar-character inventory. Thus every categorically possible
signed table is gauge equivalent to the original positive table.

For completeness, an associative signed table with these reciprocity
identities is a finite-dimensional semisimple $*$-algebra in its regular
representation. The common kernel of its regular commutators is the sum
of the one-dimensional representation spaces. Proposed scalar characters
are checked directly on every product. Their distinct irreducible
separating polynomials, whose degrees sum to the dimension of this common
kernel, prove that no scalar character is omitted. A spherical dimension
character must be nonzero on every simple and equal on dual pairs. The
remaining exact obstructions use the algebraic integrality of
$\delta/f_E$ and the necessary containment
$f_E\in\mathbb Q(d_i)$; see
\cite[Theorem~2.13 and Corollary~2.15]{OstrikPivotal}.
For a signed quotient these conditions are applied in the actual
sphericalization: its dimension is $2\delta$, its two algebra summands
are the original and signed rings, and their formal codegrees are doubled.
Consequently the same field-containment conditions apply to both summands.

The signed conclusion has a categorical consequence, rather than merely
a numerical one. Choose the lifts for which every nonzero coefficient is
positive. Their products contain only chosen lifts, so they span a fusion
subcategory of the sphericalization. Forgetting the lifts gives a tensor
equivalence with $\mathcal C$. Its spherical structure therefore descends
to $\mathcal C$; compare \cite[Section~5.1]{OstrikPivotal}.

\paragraph{Original dimension characters.}
The same complete-inventory check is performed on the fifty-nine original
rings. Actual formal codegrees, rather than uncorrected regular Casimir
eigenvalues, are used: on a matrix block belonging to $E$, the latter
eigenvalue is $(\dim E)f_E$. The scalar part is extracted first and the
remaining block sizes are checked by the dimension of the semisimple
algebra. Exact norm and field-containment certificates then reject the
inadmissible dimension characters.

For $R_2,R_{51},R_{124},R_{147}$ no spherical dimension character remains.
For all surviving characters except those of
\[
 R_{36},R_{86},R_{95},R_{135},R_{136},R_{170},
\]
the inventory supplies a field embedding and a tensor sign character
making every simple dimension positive. Galois conjugation of the
algebraic categorical data and twisting the spherical pivotal structure
by that sign character then give positive dimensions, necessarily equal
to Frobenius--Perron dimensions. Thus a Galois conjugate is pseudo-unitary;
see \cite[Section~8.4]{ENO}. This assertion does not claim that the original
category is pseudo-unitary. The six exceptions have an additional mixed
spherical dimension character and will be treated directly.

The exact replay checks fifty-nine sign records, twelve alternative signed
character records, and fifty-nine original character records, with each
record bound to the full input tensor hash. It reconstructs the sign
deductions independently. The legacy sign checker also recomputes the
twelve small alternative character calculations; their results are then
checked separately by direct homomorphism identities and completeness,
without trusting that character-discovery calculation. This full replay,
not a stored Boolean conclusion, is a dependency of the general-complex
exclusions below.

\subsection{The spherical dimension lattice}
\label{subsec:r11-dimension-lattice}

Let $d_i$ be the simple dimensions of a spherical fusion category,
$\delta=\sum_i d_i d_{i^*}$ its global dimension, and
\[
 L_d=\sum_i\mathbb Z d_i
\]
the additive lattice generated by its simple dimensions. For an
irreducible representation $E$ of its based ring with formal codegree
$f_E$, the induction theorem
\cite[Theorem~2.13]{OstrikPivotal} supplies a simple
$A_E\in\mathcal Z(\mathcal C)$ such that
\begin{equation}\label{eq:r11-central-dimension}
 \dim A_E=\frac{\delta}{f_E},\qquad
 [F(A_E):\mathbf1]=\dim E,
\end{equation}
where $F$ is the forgetful functor. Since $F$ preserves dimension,
\begin{equation}\label{eq:r11-lattice-condition}
 \frac{\delta}{f_E}\in L_d
 \qquad\text{for every }E.
\end{equation}
This condition holds for every spherical dimension character, including
characters with negative dimensions or rational global dimension.

An exact certificate for failure of
\eqref{eq:r11-lattice-condition} is a rational linear functional
$\ell$ on the character field satisfying
\[
 \ell(d_i)\in\mathbb Z\quad\text{for every }i,
 \qquad \ell(\delta/f_E)\notin\mathbb Z.
\]
The verifier evaluates these expressions directly; it does not require
Hermite normal form or a lattice search. Each irreducible character-field
factor covers its entire Galois orbit, since both the dimension condition
and the multiset of formal codegrees are Galois invariant.

Here are all thirty-four inputs rejected by the complete spherical
inventory together with these lattice witnesses. Superscripts in the
dimension types denote multiplicities.
\begin{enumerate}
\item Type $1^4,2^3,(1+\sqrt5)^4$:
\[
 2,31,51,59,69,83,107,108,122,124,147,151.
\]
The Frobenius--Perron branch has
$D=40+8\sqrt5$ and $L=\mathbb Z+\mathbb Z\sqrt5$.
The actual formal codegree $40-8\sqrt5$ requires
\[
 \frac{D}{40-8\sqrt5}=\frac{3+\sqrt5}{2}\notin L.
\]
The four inputs $2,51,124,147$ were already excluded by their complete
spherical character inventories. For the other eight, all surviving
branches are covered by exact separating functionals.

\item Type $1^4,2^3,(2+2\sqrt2)^4$:
\[
 32,62,88,110,137,138,152,154.
\]
Here $D=64+32\sqrt2$ and $L=\mathbb Z+2\mathbb Z\sqrt2$.
The formal codegree $32$ requires
\[
 \frac{D}{32}=2+\sqrt2\notin L.
\]

\item Type $1^4,2,a^2,b^2,c^2$:
\[
 34,70,89,96,115,123,139,167.
\]
Let $a$ be the largest real root of $a^3-12a-8=0$, and put
$b=a^2/2-2$ and $c=a+2$.
The numbers $1,a,b$ are rationally independent, and
\[
 L=\mathbb Z+\mathbb Za+\mathbb Zb,
 \qquad D=48+12a+12b.
\]
The formal codegree $24$ requires
\[
 \frac{D}{24}=2+\frac{a+b}{2}\notin L.
\]

\item Type $1^4,2,(1+\sqrt5)^4,(3+\sqrt5)^2$:
\[
 36,86,95,135,136,170.
\]
The Frobenius--Perron branch has $D=60+20\sqrt5$ and
$L=\mathbb Z+\mathbb Z\sqrt5$; the formal codegree $40$ requires
$(3+\sqrt5)/2\notin L$.
For each of these six rings there is also a mixed spherical character
of global dimension $\delta=40$. Its simple dimensions still generate
$\mathbb Z+\mathbb Z\sqrt5$. The actual formal codegrees
$60\pm20\sqrt5$ now require
\[
 \frac{40}{60\pm20\sqrt5}=\frac{3\mp\sqrt5}{2}
 \notin\mathbb Z+\mathbb Z\sqrt5.
\]
Thus these branches are excluded directly, without an unjustified
pseudo-unitary reduction.
\end{enumerate}

The finite replay checks forty-four separating functionals over all
remaining character branches. Four of the thirty-four inputs have no
remaining spherical branch; the other thirty require at least one lattice
witness. Combined with Section~\ref{subsec:r11-signed-bridge}, these are
exclusions of arbitrary complex categorifications.

\subsection{A two-character central-restriction obstruction}
\label{subsec:r11-two-character}

The four further rings
\[
 R_3,\quad R_{52},\quad R_{142},\quad R_{149}
\]
have dimension type $1^4,2^3,x^4$, where
$x=2+2\sqrt2$, and global Frobenius--Perron dimension
$D=64+32\sqrt2$. For these four, the complete signed and character audit
reduces existence of a complex categorification to existence of a
pseudo-unitary one. We therefore work with positive simple dimensions.

Let $\chi$ be the Galois conjugate of the Frobenius--Perron character.
It has values $1$ on the invertibles, $2$ on the two-dimensional simples,
and $2-2\sqrt2$ on the four $x$-dimensional simples. Its formal codegree is
$f_\chi=64-32\sqrt2$. A second scalar character $\varphi$ has formal
codegree $f_\varphi=16$ and takes one of the values $\pm\sqrt2$ on each
$x$-dimensional simple. In input basis order, complete witnesses for this
second character are as follows, with $q=\sqrt2$:
\begin{center}
\small
\renewcommand{\arraystretch}{1.2}
\begin{tabular}{@{}rcl@{}}
\toprule
Input & $x$-dimensional indices & $(\varphi_0,\ldots,\varphi_{10})$\\
\midrule
$3$   & $4,5,6,7$ & $(1,1,1,0,-q,q,q,-q,0,1,-2)$\\
$52$  & $4,5,7,8$ & $(1,1,1,0,-q,q,-2,q,-q,0,1)$\\
$142$ & $1,2,3,4$ & $(1,q,q,-q,-q,-2,1,0,1,1,0)$\\
$149$ & $3,4,5,6$ & $(1,1,1,q,q,-q,-q,0,-2,0,1)$\\
\bottomrule
\end{tabular}
\end{center}
Every product identity and both formal codegrees are checked exactly.

Write
\[
 \beta_\chi=\frac{D}{f_\chi}=3+2\sqrt2=1+x,
 \qquad
 \beta_\varphi=\frac{D}{f_\varphi}=4+2\sqrt2=2+x.
\]
Since both characters are one-dimensional,
\eqref{eq:r11-central-dimension} says that the restrictions of their
associated central simples contain the unit exactly once. Comparing
coefficients of $\sqrt2$ and then rational parts, using positivity, forces
\begin{equation}\label{eq:r11-forced-restrictions}
 F(A_\chi)=\mathbf1+X,
 \qquad
 F(A_\varphi)=\mathbf1+g+Y,
\end{equation}
where $X,Y$ have dimension $x$ and $g$ is a nontrivial invertible.
Indeed the irrational coefficient permits exactly one $x$-dimensional
summand, and the remaining rational weight after the prescribed unit is
respectively zero or one.

The correspondence between based-ring characters and central simples,
together with symmetry of the center's $S$-matrix, gives
\begin{equation}\label{eq:r11-reciprocity}
 \beta_\chi\,\chi(F(A_\varphi))
 =\beta_\varphi\,\varphi(F(A_\chi));
\end{equation}
see \cite[Example~2.9 and the discussion preceding Theorem~2.13]
{OstrikPivotal}. But
\[
 \chi(F(A_\varphi))=4-2\sqrt2,
 \qquad
 (3+2\sqrt2)(4-2\sqrt2)=4+2\sqrt2.
\]
Thus \eqref{eq:r11-reciprocity} forces
$\varphi(\mathbf1+X)=1$, hence $\varphi(X)=0$. This contradicts the
displayed character values on every possible $X$. It excludes all four
rings by the same argument, without enumerating full induction matrices.
Together with the preceding thirty-four, the number of distinct
general-complex exclusions in this section is thirty-eight.

\subsection{A coherent quaternionic model for $R_{179}$}
\label{subsec:r11-quaternionic-model}

We now give an existence proof for an integral input. Let $\mathcal D$ be
one of Izumi's unitary near-group categories for the quaternion group
$Q_8=\{\pm1,\pm i,\pm j,\pm k\}$. Its fusion rules are
\begin{equation}\label{eq:r11-Q8-base}
 g\rho=\rho g=\rho,
 \qquad \rho^2=\sum_{g\in Q_8}g+2\rho.
\end{equation}
Thus $\operatorname{FPdim}(\rho)=4$ and
$\operatorname{FPdim}(\mathcal D)=24$.
Existence, with the explicit categorical data used below, follows from
\cite[Theorems~1.1 and~6.1, Lemma~6.13]{Izumi}.

\paragraph{The full integral categorical tuple.}
Use the faithful quaternionic representation on $K=\mathbb C^2$ given by
\[
 I=\begin{pmatrix}i&0\\0&-i\end{pmatrix},
 \qquad J=\begin{pmatrix}0&1\\-1&0\end{pmatrix},
 \qquad K_0=IJ,
\]
so that $\pi(i)=I$, $\pi(j)=J$, and $\pi(k)=K_0$.
Let $j_0=J\circ\text{complex conjugation}$. For any third root of unity
$\zeta$, Izumi's admissible tuple is
\begin{equation}\label{eq:r11-integral-tuple}
 (K,V,U_K,j_1,j_2,\chi,l)
 =(\mathbb C^2,\pi,\pi,j_0,-\zeta j_0,\chi,0),
 \qquad
 \pi(g)\pi(h)=\chi_h(g)\pi(h)\pi(g).
\end{equation}
This is the integral tuple of Lemma~6.13, not an application of an
automorphism theorem restricted to irrational near-group categories.

Put
\begin{equation}\label{eq:r11-order-three-W}
 W=-\frac{1+I+J+K_0}{2}.
\end{equation}
Direct Gaussian-rational matrix calculations give
\[
 W^*W=1,\qquad W^3=1,\qquad \operatorname{tr}(W)=-1,
 \qquad Wj_0=j_0W.
\]
Conjugation by $W$ cyclically permutes $I,J,K_0$. Hence, for the order-three
group automorphism $\phi:i\mapsto j\mapsto k\mapsto i$,
\[
 W\pi(g)W^*=\pi(\phi(g))\qquad(g\in Q_8).
\]
These equations preserve both representations $V$ and $U_K$ in
\eqref{eq:r11-integral-tuple}, both antiunitaries $j_1,j_2$, and the
commutator bicharacter $\chi$ after relabeling by $\phi$. Since $l=0$, its
compatibility is automatic. This verifies every component of the
equivalence-of-tuples condition in \cite[Definition~4.8]{Izumi}.

\paragraph{Categorical action and coherence.}
For clarity, existence of the action can be read directly from the
reconstruction, not merely from a permutation of fusion rules.
The generating isometries are
$S_g\in\operatorname{Hom}(g,\rho^2)$ and
$T(v)\in\operatorname{Hom}(\rho,\rho^2)$, $v\in K$.
Define the Cuntz-algebra automorphism
\begin{equation}\label{eq:r11-Cuntz-action}
 \gamma(S_g)=S_{\phi(g)},\qquad
 \gamma(T(v))=T(Wv).
\end{equation}
The reconstruction equations \cite[(3.10)--(3.12)]{Izumi} are covariant
under this assignment. In (3.10), $\sum_gS_g$ and the contraction
$\sum_a T_a j_1(T_a)$ are preserved. Equation (3.11) is covariant because
the pointed representations intertwine. Each term in (3.12) is covariant
because $W$ intertwines $V,U_K$, commutes with $j_1,j_2$, and preserves
$l=0$. Therefore $\gamma$ conjugates the pointed endomorphisms by $\phi$
and commutes with $\rho$. The reconstruction theorem
\cite[Theorem~4.9]{Izumi} supplies the resulting unitary tensor
autoequivalence of $\mathcal D$.

All morphisms are generated by these isometries and their adjoints
\cite[Remark~3.6]{Izumi}. Since $\phi^3=1$ and $W^3=1$, the cube of
\eqref{eq:r11-Cuntz-action} is the identity on the Cuntz generators,
hence on the reconstructed category. Identity compositors consequently
give a coherent, strict $C_3$ action. No use is made of Izumi's Section~13
automorphism results, which have a different, irrational scope.

\paragraph{Equivariant simples and products.}
The unitary equivariantization
$\mathcal C=\mathcal D^{C_3}$ has dimension $72$.
The fixed invertibles $1,-1$ each have three equivariant structures,
giving six invertibles $L_{(g,a)}$, $g\in\{1,-1\}$,
$a\in\mathbb Z/3$. The free orbits
$\{i,j,k\}$ and $\{-i,-j,-k\}$ give two objects $U,V$ of dimension three.
The fixed object $\rho$ has three lifts $P_a$ of dimension four.
The rank and dimension type are therefore $11$ and $1^6,3^2,4^3$.

The pointed part is computed by equivariant homogeneous fibers. Explicitly,
\begin{align*}
 L_{(g,a)}L_{(h,b)}&=L_{(gh,a+b)},\\
 L_{(1,a)}U&=U L_{(1,a)}=U,
 &L_{(1,a)}V&=V L_{(1,a)}=V,\\
 L_{(-1,a)}U&=U L_{(-1,a)}=V,
 &L_{(-1,a)}V&=V L_{(-1,a)}=U,\\
 U^2=V^2&=\sum_{a\in\mathbb Z/3}L_{(-1,a)}+U+V,
 &UV=VU&=\sum_{a\in\mathbb Z/3}L_{(1,a)}+U+V.
\end{align*}
The remaining products are
\begin{align}
 L_{(g,a)}P_b=P_bL_{(g,a)}&=P_{a+b},\label{eq:r11-fixed-rho}\\
 UP_b=P_bU=VP_b=P_bV&=P_0+P_1+P_2,\label{eq:r11-orbit-rho}\\
 P_aP_b&=L_{(1,a+b)}+L_{(-1,a+b)}+U+V\notag\\
 &\hspace{12mm}+P_{a+b+1}+P_{a+b+2}.
 \label{eq:r11-rho-square}
\end{align}
Indices are read modulo three. In \eqref{eq:r11-rho-square}, the fixed
pointed channels have phase one because $S_1,S_{-1}$ are fixed by
$\gamma$. The action on $\operatorname{Hom}(\rho,\rho^2)$ is $W$, whose
eigencharacters are the two nontrivial characters of $C_3$. An inverse
convention for the action replaces $W$ by $W^{-1}$ and leaves this
multiset unchanged. In the right absorption channels of
\eqref{eq:r11-fixed-rho}, no additional phase occurs: $\pi(\pm1)$ commutes
with $W$. Free three-cycles contribute each character once, yielding
\eqref{eq:r11-orbit-rho}. Finally $P_a^*=P_{-a}$, so exactly one
four-dimensional simple is self-dual.

In the model basis
\[
 L_{(1,0)},L_{(1,1)},L_{(1,2)},
 L_{(-1,0)},L_{(-1,1)},L_{(-1,2)},U,V,P_0,P_1,P_2,
\]
the model-to-input bijection is
\begin{equation}\label{eq:r11-179-bijection}
 (0,7,8,9,1,2,5,6,10,4,3).
\end{equation}
Every one of the $11^3=1331$ multiplication coefficients agrees under
this bijection with $R_{179}$. Its canonical tensor SHA-256 is
\begin{center}\small
\texttt{a9492a728c3c2e5d59d92d723a5349782a}\\[-1mm]
\texttt{962bc04ead4aacd0d7756e5d7c95f2}.
\end{center}
All three cube-root choices in \eqref{eq:r11-integral-tuple} provide
this based tensor. One such realization suffices for the existence claim.

\subsection{Finite verification and dependency order}
\label{subsec:r11-nonintegral-replay}

All following paths are relative to \texttt{work/rank11/nonintegral/}.
Run the four accepted entrypoints in this order, from the project root:
\begin{enumerate}
\item \nolinkurl{verify_reduction.py}: complete signed and scalar replay;
59 sign records, 12 alternative signed character records, and
59 original character records.
\item \nolinkurl{verify_spherical_lattice.py}: 34 excluded inputs,
including the four already rejected scalar inventories, and
44 exact separating functionals.
\item \nolinkurl{verify_two_character.py}: the four explicit character
pairs, formal codegrees, forced restrictions, and their contradictions.
Its rational arithmetic is implemented directly in $\mathbb Q(\sqrt2)$.
\item \nolinkurl{integral179/verify.py}: an independent reconstruction of
the equivariant fusion tensor by
\[
 \dim H^{C_3}=\frac13\sum_{t=0}^2\operatorname{Tr}(t\mid H)
\]
on every underlying near-group fusion space. The exceptional
two-dimensional space has character $(2,-1,-1)$; the other channels
are the one-dimensional absorption or $S_g$ spaces, with their orbit
permutations. This verifier does not import the model constructor.
It checks the exact action matrix, the source hash, all $1331$ tensor
coefficients, and all $11^4=14641$ associativity identities.
\end{enumerate}

The first replay imports the project's prior oriented-sign and exact
character verifiers. The lattice replay imports the shared exact-number
field helpers. The two-character and quaternionic model verifiers use
only the Python standard library. Critical conditions use explicit
exceptions and remain active under \texttt{python -O}. In particular,
the second and third checks alone are not certificates of a
general-complex exclusion: they require the full first replay, not merely
its summary flags. No timeout, failed model search, numerical tolerance,
or incomplete finite domain is used in any accepted conclusion here.
